% curva-rss.tex — Figura obligatoria 3 (R-30): curva de memoria durante la ventana.
%
% Dos series por panel, y la diferencia entre ellas es el hallazgo F6-A3:
%   anon        = lo que mide la columna `max_rss_mb` del CSV;
%   working_set = memory.current − inactive_file, que SÍ incluye la memoria
%                 compartida (`shared_buffers`) de PostgreSQL.
% Los `shared_buffers` son `shmem`, no `anon`, así que con `max_rss_mb` PostgreSQL
% sale artificialmente eficiente (SPEC §4 y A11). M3 se interpreta con working_set.
%
% Fuente: la serie de 1 s del cgroup que el runner guarda en results/raw/
% (serie_cgroup_1s), es decir las MISMAS lecturas de las que sale la columna.

\newcommand{\rssPanel}[4]{%
  % #1 archivo .dat · #2 título · #3 color · #4 "leyenda" para mostrarla
  \begin{tikzpicture}
    \begin{axis}[
        title={#2}, title style={font=\small},
        xlabel={Segundo de la ventana de medición}, ylabel={Memoria (MiB)},
        xmin=0, xmax=300, ymin=0,
        width=0.47\linewidth, height=5cm,
        legend pos=south east, legend cell align=left,
      ]
      \addplot[color=#3, thick, no marks] table [x=t_s, y=working_set_mib] {#1};
      \addplot[color=#3!55, thick, densely dashed, no marks]
        table [x=t_s, y=anon_mib] {#1};
      \addplot[color=gray!60, thin, dotted, no marks] table [x=t_s, y=shmem_mib] {#1};
      \ifthenelse{\equal{#4}{leyenda}}%
        {\legend{\texttt{working\_set}, \texttt{anon}, \texttt{shmem}}}{}%
    \end{axis}
  \end{tikzpicture}%
}

\begin{figure}[htbp]
  \centering
  \rssPanel{rss_mongodb_A_Q2_c50.dat}{MongoDB · perfil A · Q2 · c=50}{blue!60!black}{leyenda}\hfill
  \rssPanel{rss_postgresql_A_Q2_c50.dat}{PostgreSQL · perfil A · Q2 · c=50}{red!60!black}{}

  \rssPanel{rss_mongodb_C_Q2_c50.dat}{MongoDB · perfil C · Q2 · c=50}{blue!60!black}{}\hfill
  \rssPanel{rss_postgresql_C_Q2_c50.dat}{PostgreSQL · perfil C · Q2 · c=50}{red!60!black}{}

  \caption[Curva de memoria durante la ventana]{%
    \textbf{Memoria del contenedor durante la ventana, muestreada cada segundo.}
    La línea discontinua (\texttt{anon}) es exactamente lo que resume la columna
    \texttt{max\_rss\_mb}; la continua (\texttt{working\_set}) incluye además la
    memoria compartida. \textbf{En PostgreSQL la distancia entre las dos es
    \texttt{shared\_buffers}}: son \texttt{shmem}, no \texttt{anon}, así que
    \texttt{max\_rss\_mb} no cuenta su búfer principal y lo hace parecer
    artificialmente eficiente frente a MongoDB, cuyo caché de WiredTiger sí es
    \texttt{anon} (F6-A3, SPEC §4 y A11). Por eso M3 se interpreta con
    \texttt{working\_set}. Fuente: serie de 1 s del cgroup en
    \texttt{results/raw/}, las mismas lecturas de las que sale la columna.}
  \label{fig:curva-rss}
\end{figure}
